% AC-G112-twist: new finite-character classifier construction.
% No mixed Selmer or integral model-length conclusion.
\subsection{Cubic endpoint twists in the common trace classifier}
\label{sec:cubic-twisted-block-trace-data}

We retain the exact source notions \texttt{Af.TrData} and
\texttt{Af.TrData.Classifies} from \texttt{Classifies.lean}.
An \texttt{Af.TrData}(Q,e) has a finite carrier, a complex trace
function, a permutation for every complex field automorphism, the
bound three, and integrality after multiplication by \(Q^e\).
Its classifier is injective up to isomorphism of representations and
evaluates the actual trace at the specified group element.  No
classifier of a representation is asserted to be canonical.

Fix a group \(\Gamma\), an element \(g\), and a fixed quotient
\(\theta:\Gamma\to C_3\).  The group and quotient are independent of
the coefficient prime.  Let \(Q,e\) be positive integers.  Suppose
one is given a finite set \(S_0\) and a map \(t_0:S_0\to\mathbb C\)
such that
\begin{align*}
 &\text{for every }\sigma\in\operatorname{Aut}(\mathbb C),\text{ some
 permutation }u_\sigma\text{ satisfies }
          t_0(u_\sigma s)=\sigma(t_0(s));\\
 &Q^e t_0(s)\text{ is an algebraic integer for every }s.
                                                        \tag{G112T.1}
\end{align*}
For every admitted coefficient prime \(\ell\) and isomorphism
\(\iota:\overline{\mathbb Q}_\ell\simeq\mathbb C\), assume these data
classify the admitted semisimple representations
\[
                  R_0(V)=V\oplus\mathbf1
                                                        \tag{G112T.2}
\]
by the precise injection and trace-evaluation conditions of the
source \texttt{Classifies}.  The old carrier and function are selected
before \(\ell\).  These hypotheses hold if the old classifier is
either \texttt{Win.TrData} or \texttt{Af.TrData}; its old size field
will not be used.  This distinction avoids confusing the weaker
\(\sqrt Q+2\) bound in the former namespace with the bound three in
the latter.

Retain also the following purity premise for every admitted block:
\[
 \bigl|\sigma(\iota\operatorname{tr}V(g))\bigr|\le2
 \quad\text{for every }\sigma\in\operatorname{Aut}(\mathbb C).
                                                        \tag{G112T.3}
\]
The order of the middle and endpoint coordinates in (G112T.2) is
irrelevant to its isomorphism class.  No continuity hypothesis is
needed for the finite-data construction; in the arithmetic application
all representations are the actual continuous ones.

\begin{proposition}[AC-G112-twist: Cubic endpoint-twist trace classifier]
\label{prop:G112-twist}
Under (G112T.1)--(G112T.3) there is an
\(\texttt{Af.TrData}(Q,e)\) object \(D_\chi\), selected before the
coefficient prime, whose carrier has cardinality at most \(3\#S_0\).
It classifies every admitted semisimple representation isomorphic to
\[
 R(V,\tau)=\bigl(V\otimes(\tau\circ\theta)\bigr)\oplus\mathbf1,
       \qquad \tau:C_3\to\overline{\mathbb Q}_\ell^\times,
                                                        \tag{G112T.4}
\]
where \(V\) is admitted in (G112T.2).  In particular this gives the
actual cubic Type II trace classifier whenever the untwisted old
block theorem and its purity/counting inputs have been supplied.
It does not change the literal source definition \texttt{Adm2B}.
\end{proposition}

\begin{proof}
Take the finite subset
\[
 S^\circ=\{s\in S_0:
       |\sigma(t_0(s))-1|\le2\text{ for every }
                  \sigma\in\operatorname{Aut}(\mathbb C)\}.
                                                        \tag{G112T.5}
\]
Any permutation \(u_\sigma\) in (G112T.1) preserves this subset:
for every automorphism \(\tau'\),
\(\tau'(t_0(u_\sigma s))=(\tau'\circ\sigma)(t_0(s))\), and
composition with \(\sigma\) permutes all automorphisms.  The same
argument, or the equality with \(\sigma^{-1}\) applied to
\(t_0(u_\sigma s)\), gives the converse.  Thus the restriction is a
permutation.  The old trace identity and (G112T.3) show that every
actual classifier label of (G112T.2) lies in \(S^\circ\).

Put
\[
 Z=\operatorname{Hom}(C_3,\mu_3(\mathbb C)),\qquad
 D_\chi.S=S^\circ\times Z.
                                                        \tag{G112T.6}
\]
The set \(Z\) has three elements.  For \(\xi\in Z\), write
\(z=\xi(\theta(g))\), and define
\[
                   t_\chi(s,\xi)=1+z\bigl(t_0(s)-1\bigr).
                                                        \tag{G112T.7}
\]
For a complex field automorphism \(\sigma\), the permutation
\((s,\xi)\mapsto(u_\sigma s,\sigma\circ\xi)\) carries this trace
to its \(\sigma\)-conjugate.  Since \(|z|=1\), (G112T.5) with the
identity automorphism gives \(|t_\chi(s,\xi)|\le3\).  Moreover
\[
 Q^e t_\chi(s,\xi)=zQ^e t_0(s)+(1-z)Q^e
                                                        \tag{G112T.8}
\]
is an algebraic integer by (G112T.1) and the integrality of roots of
unity.  These prove every field of the literal
\texttt{Af.TrData}(Q,e), including its stronger size field.
The carrier is finite of cardinality at most \(3\#S_0\), and its
definition involves no coefficient prime or coefficient-field embedding.

For each admitted representation choose one isomorphism and one
witness \((V,\tau)\) as in (G112T.4).  Take its old classifier label
\(s\in S^\circ\) and assign the new label
\((s,\iota\circ\tau)\).  The second coordinate is indeed an element
of \(Z\), since the image of \(\tau\) consists of third roots of
unity.  The trace of (G112T.4) at \(g\) is (G112T.7), transported
through \(\iota^{-1}\).  Trace invariance under the chosen
isomorphism gives the exact source evaluation condition.

If two chosen labels coincide, the old classifier gives
\(V\oplus\mathbf1\simeq V'\oplus\mathbf1\).  Semisimplicity and
cancellation of the same simple summand give \(V\simeq V'\).
This can also be seen from uniqueness of the simple-constituent
multiplicities in a finite-dimensional semisimple representation.
Equality of the second coordinates and injectivity of \(\iota\)
give equality of the whole characters \(\tau=\tau'\).  Twisting
the isomorphism of blocks and adjoining the middle line give an
isomorphism of the two representations (G112T.4).  Hence the map is
injective up to the precise source \texttt{RepIso}, as required.

It is essential that (G112T.6) retains the whole character.  At a
test place with \(\theta(g)=1\), distinct cubic twists have the same
value at \(g\).  Indexing only by that value would not justify the
classifier injection.  No such identification is made here.
\end{proof}

\paragraph{Exact old modular interface and its scope.}
Use \(\Gamma=G_K\), the fixed Kummer quotient of
\(K(\sqrt[3]{269})/K\), and geometric Frobenius at a test place
\(w\mid q\).  AC-G68-block/point at 269, and their already licensed
analogues at other coefficient primes, give
\[
       \rho\simeq (V_f|_{G_K}\otimes\chi^j)\oplus\mathbf1.
                                                        \tag{G112T.9}
\]
The rational representation \(V_f\) is continuous, irreducible,
has determinant \(\eta_K\), and satisfies the old point theorem;
its restriction to \(G_K\) is the untwisted endpoint block.
Under that
theorem's unchanged old automorphy input, including its \(e_O=1\)
condition where invoked, its \texttt{CohInp}, and its literal
\texttt{GD}, the full source
\texttt{adm2B\_rhoK\_of\_typeII\_of\_modular\_extensions} applies
to the \emph{untwisted} point.  It gives the old
\texttt{Adm2B} representation \(V_f|_{G_K}\oplus\mathbf1\), at
the original level \(M_0=N|D_K|\).  No assertion is made that
(G112T.9) itself has this old admissibility property.

Retain exactly \texttt{CountQAlt} and \texttt{NewAdmAlt}, as used
in the complete source proof of
\texttt{frobenius\_trace\_count\_bound\_of\_rat\_eigensystem\_count\_alt}.
Alternatively retain its displayed hypotheses
\texttt{Hypotheses'}, \texttt{LevelPNAlt}, \texttt{AdmResNAlt},
and \texttt{NewAdmAlt}, which supply \texttt{CountQAlt} in the
source.  These inputs give, for \(Q=q^{h}=Nw\), one old carrier
\(S_0\) with
\[
              \#S_0\le c_2\{N|D_K|(e+1)\}^{A_2},
                                                        \tag{G112T.10}
\]
chosen before \(\ell\ne2,q\).  The source proof includes the
residue-degree power \(h\), inverse Frobenius, the middle line,
and the precise trace equality.  We use its finite carrier,
conjugation, integrality and classifier fields, whether the source
wrapper is in the \texttt{Win} or \texttt{Af} namespace.  Its size
field is not used in AC-G112-twist.

The retained modular purity input verifies (G112T.3).  In the exact
source normalization, a newform \(f\) of weight \(2e+1\) is attached
to \(V_f c_\ell^{-e}\).  At a good rational prime, its two
geometric Frobenius eigenvalues have complex absolute value \(q^e\)
under every embedding of its algebraic coefficient field.  Since
\(c_\ell(\mathrm{Frob}^{\rm geom}_q)=q^{-1}\), the eigenvalues of
\(V_f\) have absolute value one.  Taking the \(h\)-th power for
restriction to \(K\), or inverse Frobenius, preserves this assertion.
Every complex automorphism restricts to an embedding of the algebraic
coefficient field, so the assertion holds for all automorphisms in
(G112T.3), not just for one chosen embedding.  This is the usual
purity theorem for the admitted modular blocks, retained as an
explicit imported arithmetic input; no new proof of it is asserted.

There is no compatibility requirement selecting one complex primitive
cube root simultaneously for all \(\ell\).  For each actual
coefficient realization of \(\chi^j\), its composite with \(\iota\)
is one of the three elements of \(Z\).  The common carrier contains
all of them.  Finite coefficient enlargements preserve the block
isomorphisms and the character labels.  In particular the old
coefficient uniformity is not lost by adjoining cubic character values.

\paragraph{The common integer after both classifier replacements.}
Let \(K=\mathbb Q(\sqrt{-3})\), \(d=|D_K|=3\),
\(Y=N(e+1)\ge1\), and let \(q\ge7\) be a rational test prime
avoiding \(269Nd\).  For \(w\mid q\), write \(Q=q^h\) with
\(1\le h\le2\).  AC-G112-count gives an
\texttt{Af.TrData} object of cardinality at most \(C_1Y^{A_1}\)
for the stated Type I domain.  Equations (G112T.6) and (G112T.10)
give one of cardinality at most \(3c_2d^{A_2}Y^{A_2}\) for the
stated cubic Type II domain.  Their literal \texttt{Af.TrData.sum}
classifies the union, with cardinality at most
\[
       (C_1+3c_2d^{A_2})Y^{A_1+A_2}.
                                                        \tag{G112T.11}
\]
The source \texttt{Classifies.sum} proof applies without alteration;
it selects one summand for an object in both domains and uses the
disjoint-union labels for the injection.

For any finite Galois-stable family of roots of unity \(z(o)\),
the exact theorem
\texttt{UpperTheory.frobenius\_trace\_norm\_nonzero\_log\_bound}
therefore gives one integer \(\mathcal N\ne0\), before \(\ell\),
with all of its stated trace-difference and \(\ell\)-adic product
bounds on the admitted union.  Explicitly, for each admitted
\(\ell\ne q\), each \(\iota\), and finite sets \(X_o\) of such
representations which are pairwise nonisomorphic within each \(X_o\),
all differences
\(\operatorname{tr}\rho(g)-\iota^{-1}(\operatorname{eis3}(Q,e,z(o)))\)
are nonzero and
\[
 |\mathcal N|_\ell\le
 \prod_o\prod_{\rho\in X_o}
 \bigl|\operatorname{tr}\rho(g)-
       \iota^{-1}(\operatorname{eis3}(Q,e,z(o)))\bigr|_\ell.
                                                        \tag{G112T.11a}
\]
Its logarithmic bound is at most
\[
 5(C_1+3c_2d^{A_2})\,\#\mathrm{Orb}\,
              Y^{A_1+A_2+1}\log q.
                                                        \tag{G112T.12}
\]
Indeed \(\log Q\le2\log q\), \(\log q\ge1\), and \(e+1\le Y\)
give \(2e\log Q+1\le5Y\log q\); substitute (G112T.11) into the
literal norm theorem.  This is exactly the scalar counting algebra
in \texttt{KeepNumU}, now fed by the two verified \texttt{Af} data
objects instead of an unproved identification of namespaces or of
cubic Type II with old \texttt{Adm2B}.  The test prime condition
\(\ell\ne q\) supplies \(\ell\nmid Q\).

The finite-data norm theorem above is valid at any test prime in its
stated domain.  Identifying its \(\operatorname{eis3}(Q,e,z)\) with
the \emph{actual cubic seed trace} retains the separately proved
controlled inert test-prime choice of AC-K29, for which
\(\chi(\mathrm{Frob}_w)=1\).  At a split test prime with nontrivial
cubic value, the same cubic scalar multiplies both endpoint terms;
it cannot in general be absorbed into the reciprocal pair
\(z,z^{-1}\).  No such seed-trace identification at arbitrary split
primes is asserted by this construction.

\paragraph{Remaining boundary.}
This removes the finite cubic endpoint-twist classifier proviso on
the displayed old modular/counting domain.  All old automorphy,
coefficient, purity and counting hypotheses remain visible.  The
result gives neither mixed Bloch--Kato vanishing at 269 nor the
integral model-length comparison that would be needed to bound the
original deformation quantity.  It does not create models at the
omitted coefficient primes or prove the final numerical budget.
The constants are uniform in \(e,N,\ell\) only to the same extent
as the retained old counting constants; the new factors three and
269 are fixed.

\paragraph{Geometric dictionary.}
The normalization of \(\operatorname{Spec}\mathcal O_K[1/269]\)
in \(K(\sqrt[3]{269})\) is the fixed finite etale cubic cover.
In the already audited finite-cover dictionary of SGA~1,
Expose~V, sections~4 and~7, its three characters give the three
rank-one Artin local systems used in (G112T.6).  On a common open
subset where the endpoint local system is defined, (G112T.4)
tensors only that endpoint system by the chosen Artin system and
retains the constant middle line.  The finite trace carrier is a
classifier for these representations.  It constructs no new motive
and implies no Selmer vanishing.  This paragraph uses only the
previously audited dictionary; it makes no new source-scan claim.
